% GENERATED FILE. Edit canonical lessons, then run npm run generate:books.
% canonical-source-hash: fnv1a64:6e4596e3130d1b1f
% canonical-lessons: BN-C180-rojgar-kora, BN-C180-dkhashona-kora, BN-C180-pala, BN-C180-jol-deoya, BN-C180-chhapa

\chapter{To earn, To look after, To rear, To water, To print}
\label{ch:bn-to-earn-to-look-after-to-rear-to-water-to-print}
\hlchaptermodality{180}

\begin{chapteropening}
\textbf{By the end of this chapter:} \emph{I can say to earn, to look after, to rear, to water and to print.}

Complete the last lesson of chapter 180: I can say to earn, to look after, to rear, to water and to print.
\end{chapteropening}

\section[\texorpdfstring{\bn{রোজগার} \bn{করা} (rojgār kôrā)}{রোজগার করা (rojgār kôrā)}]{\bn{রোজগার} \bn{করা} (rojgār kôrā) — to earn}

\label{lesson:BN-C180-rojgar-kora}



\begin{quote}
Before the new one: say the Bengali for to look, then the Bengali for to pick.
\end{quote}

\subsection*{You'll want to know: \bn{রোজগার} \bn{করা}}
\textbf{\bn{রোজগার} \bn{করা}} — \emph{rojgār kôrā} — \textquotedblleft{}to earn\textquotedblright{}.

\begin{cousinweb}
\textbf{\bn{রোজগার} \bn{করা}} (\emph{rojgār kôrā}) is \emph{rojgār}, \textquotedblleft{}earnings\textquotedblright{}, usually traced to Persian \emph{rōzgār}, \textquotedblleft{}time, fortune; livelihood\textquotedblright{}, + \emph{kôrā}, \textquotedblleft{}to do\textquotedblright{}, from Sanskrit \emph{karoti}, \textquotedblleft{}does\textquotedblright{} (CDIAL 2814).
\end{cousinweb}

\begin{grammarlens}[title={What you've built}]
One more word for this chapter.
\end{grammarlens}

\subsection*{Your turn}
\begin{itemize}
\raggedright
  \item \emph{Say it:} \emph{rojgār kôrā}
  \item \emph{Say it:} \emph{rojgār kôrā}, once more
  \item \emph{Say it:} say \emph{rojgār kôrā}
  \item \emph{Recall:} say the Bengali for to look, then the Bengali for to pick, then say \emph{rojgār kôrā} again
\end{itemize}

\begin{tcolorbox}[breakable,colback=teal!4,colframe=teal!35!black,title={Before you move on}]
What does \bn{রোজগার} \bn{করা} mean? (\textquotedblleft{}To earn\textquotedblright{}.) Say it once more.
\end{tcolorbox}
\section[\texorpdfstring{\bn{দেখাশোনা} \bn{করা} (dækhāshonā kôrā)}{দেখাশোনা করা (dækhāshonā kôrā)}]{\bn{দেখাশোনা} \bn{করা} (dækhāshonā kôrā) — to look after}

\label{lesson:BN-C180-dkhashona-kora}



\begin{quote}
Before the new one: say the Bengali for to pick, then the Bengali for to earn.
\end{quote}

\subsection*{You'll want to know: \bn{দেখাশোনা} \bn{করা}}
\textbf{\bn{দেখাশোনা} \bn{করা}} — \emph{dækhāshonā kôrā} — \textquotedblleft{}to look after\textquotedblright{}.

\begin{cousinweb}
\textbf{\bn{দেখাশোনা} \bn{করা}} (\emph{dækhāshonā kôrā}) is \emph{dækhā}, \textquotedblleft{}seeing\textquotedblright{}, which CDIAL lists under \emph{*dṛkṣati}, \textquotedblleft{}sees\textquotedblright{} (CDIAL 6507), + \emph{shonā}, \textquotedblleft{}hearing\textquotedblright{}, from Sanskrit \emph{śṛṇoti}, \textquotedblleft{}hears\textquotedblright{} (CDIAL 12598), \textquotedblleft{}seeing and hearing\textquotedblright{}, + \emph{kôrā}, \textquotedblleft{}to do\textquotedblright{}, from Sanskrit \emph{karoti}, \textquotedblleft{}does\textquotedblright{} (CDIAL 2814).
\end{cousinweb}

\begin{grammarlens}[title={What you've built}]
One more word for this chapter.
\end{grammarlens}

\subsection*{Your turn}
\begin{itemize}
\raggedright
  \item \emph{Say it:} \emph{dækhāshonā kôrā}
  \item \emph{Say it:} \emph{dækhāshonā kôrā}, once more
  \item \emph{Say it:} say \emph{dækhāshonā kôrā}
  \item \emph{Recall:} say the Bengali for to pick, then the Bengali for to earn, then say \emph{dækhāshonā kôrā} again
\end{itemize}

\begin{tcolorbox}[breakable,colback=teal!4,colframe=teal!35!black,title={Before you move on}]
What does \bn{দেখাশোনা} \bn{করা} mean? (\textquotedblleft{}To look after\textquotedblright{}.) Say it once more.
\end{tcolorbox}
\section[\texorpdfstring{\bn{পালা} (pālā)}{পালা (pālā)}]{\bn{পালা} (pālā) — to rear, to keep (animals)}

\label{lesson:BN-C180-pala}



\begin{quote}
Before the new one: say the Bengali for to earn, then the Bengali for to look after.
\end{quote}

\subsection*{You'll want to know: \bn{পালা}}
\textbf{\bn{পালা}} — \emph{pālā} — \textquotedblleft{}to rear, to keep (animals)\textquotedblright{}.

\begin{cousinweb}
\textbf{\bn{পালা}} (\emph{pālā}) is a tadbhava, a word inherited from Sanskrit and reshaped by sound change. It comes from \emph{pālayati}, \textquotedblleft{}guards, protects\textquotedblright{} (CDIAL 8129). CDIAL lists Bengali \emph{pālā}, \textquotedblleft{}to rear, keep\textquotedblright{}.
\end{cousinweb}

\begin{grammarlens}[title={What you've built}]
One more word for this chapter.
\end{grammarlens}

\subsection*{Your turn}
\begin{itemize}
\raggedright
  \item \emph{Say it:} \emph{pālā}
  \item \emph{Say it:} \emph{pālā}, once more
  \item \emph{Say it:} say \emph{pālā}
  \item \emph{Recall:} say the Bengali for to earn, then the Bengali for to look after, then say \emph{pālā} again
\end{itemize}

\begin{tcolorbox}[breakable,colback=teal!4,colframe=teal!35!black,title={Before you move on}]
What does \bn{পালা} mean? (\textquotedblleft{}To rear, to keep (animals)\textquotedblright{}.) Say it once more.
\end{tcolorbox}
\section[\texorpdfstring{\bn{জল} \bn{দেওয়া} (jôl deoyā)}{জল দেওয়া (jôl deoyā)}]{\bn{জল} \bn{দেওয়া} (jôl deoyā) — to water}

\label{lesson:BN-C180-jol-deoya}



\begin{quote}
Before the new one: say the Bengali for to look after, then the Bengali for to rear.
\end{quote}

\subsection*{You'll want to know: \bn{জল} \bn{দেওয়া}}
\textbf{\bn{জল} \bn{দেওয়া}} — \emph{jôl deoyā} — \textquotedblleft{}to water\textquotedblright{}.

\begin{cousinweb}
\textbf{\bn{জল} \bn{দেওয়া}} (\emph{jôl deoyā}) is \emph{jôl}, \textquotedblleft{}water\textquotedblright{}, from Sanskrit \emph{jala} (CDIAL 5155), + \emph{deoyā}, \textquotedblleft{}to give\textquotedblright{}, from Sanskrit \emph{dadāti}, \textquotedblleft{}gives\textquotedblright{} (CDIAL 6141): \textquotedblleft{}to give water\textquotedblright{}.
\end{cousinweb}

\begin{grammarlens}[title={What you've built}]
One more word for this chapter.
\end{grammarlens}

\subsection*{Your turn}
\begin{itemize}
\raggedright
  \item \emph{Say it:} \emph{jôl deoyā}
  \item \emph{Say it:} \emph{jôl deoyā}, once more
  \item \emph{Say it:} say \emph{jôl deoyā}
  \item \emph{Recall:} say the Bengali for to look after, then the Bengali for to rear, then say \emph{jôl deoyā} again
\end{itemize}

\begin{tcolorbox}[breakable,colback=teal!4,colframe=teal!35!black,title={Before you move on}]
What does \bn{জল} \bn{দেওয়া} mean? (\textquotedblleft{}To water\textquotedblright{}.) Say it once more.
\end{tcolorbox}
\section[\texorpdfstring{\bn{ছাপা} (chhāpā)}{ছাপা (chhāpā)}]{\bn{ছাপা} (chhāpā) — to print}

\label{lesson:BN-C180-chhapa}



\begin{quote}
Before the new one: say the Bengali for to rear, then the Bengali for to water.
\end{quote}

\subsection*{You'll want to know: \bn{ছাপা}}
\textbf{\bn{ছাপা}} — \emph{chhāpā} — \textquotedblleft{}to print\textquotedblright{}.

\begin{cousinweb}
\textbf{\bn{ছাপা}} (\emph{chhāpā}) is a tadbhava, a word inherited and reshaped by sound change. It comes from \emph{*chapp}, \textquotedblleft{}press, cover, hide\textquotedblright{}, a form CDIAL reconstructs from its descendants (CDIAL 4994). CDIAL lists Bengali \emph{chāpā}, \textquotedblleft{}to be concealed, to print\textquotedblright{}.
\end{cousinweb}

\begin{grammarlens}[title={What you've built}]
One more word for this chapter.
\end{grammarlens}

\subsection*{Your turn}
\begin{itemize}
\raggedright
  \item \emph{Say it:} \emph{chhāpā}
  \item \emph{Say it:} \emph{chhāpā}, once more
  \item \emph{Say it:} say \emph{chhāpā}
  \item \emph{Recall:} say the Bengali for to rear, then the Bengali for to water, then say \emph{chhāpā} again
\end{itemize}

\begin{tcolorbox}[breakable,colback=teal!4,colframe=teal!35!black,title={Before you move on}]
What does \bn{ছাপা} mean? (\textquotedblleft{}To print\textquotedblright{}.) Say it once more.
\end{tcolorbox}
